% !PRINTED 388
\weberpage{388}

%<<CONT
und es genügt also, den ganzen Kettenbruch durch die Periode zu bezeichnen, etwa so:

\weq{10}{\omega = [\alpha_0, \alpha_1 \ldots \alpha_{\nu-1}].}

Dieser Satz ist offenbar gleichbedeutend damit, dass die Schlusszahl $\omega_\nu$ des Kettenbruches (9) mit $\omega$ selbst identisch ist, und darin liegt auch der Beweis der Behauptung. Denn die Schlusszahlen $\omega_\nu$ gehören als äquivalente Zahlen alle zur selben Discriminante, und daher ist nach 4. die Zahl aller möglichen Schlusszahlen nur eine endliche. In der Reihe der Schlusszahlen $\omega, \omega_1, \omega_2 \ldots$ muss daher einmal eine schon dagewesene zum zweitenmal auftreten. Ist $\omega_\nu$ die erste, die zum zweitenmal auftritt, so muss $\omega_\nu = \omega$ sein, da sonst nach 3. \S.\,124 auch $\omega_{\nu-1}$ zum zweitenmal auftreten würde.

Es lassen sich also die sämmtlichen, zu einer Discriminante gehörigen reducirten Zahlen in Perioden anordnen:

\weq{11}{\omega, \omega_1, \omega_2 \ldots \omega_{\nu-1},}

so dass, wenn
\[
\alpha_0, \alpha_1, \alpha_2 \ldots \alpha_{\nu-1}
\]
die grössten darin enthaltenen ganzen Zahlen sind
\[
\begin{array}{l}
\omega = [\alpha_0, \alpha_1 \ldots \alpha_{\nu-1}]\\
\omega_1 = [\alpha_1, \alpha_2 \ldots \alpha_0]\\
\cdots\cdots\cdots\cdots\cdots\cdots\cdots\cdots
\end{array}
\]
ist, womit die Kettenbruchentwickelung für jede von diesen Zahlen gegeben ist.

Ist durch \gesperrt{eine} Periode das ganze System der reducirten Zahlen noch nicht erschöpft, so bildet man eine zweite Periode u. s. f. Da durch eine Zahl $\omega$ sowohl die in der Periode vorhergehende, wie die nachfolgende völlig bestimmt ist, so enthalten zwei verschiedene Perioden niemals eine gemeinschaftliche Zahl.

Von diesen Perioden gilt nun der Satz:

\gesperrt{6. Reducirte Zahlen aus derselben Periode sind äquivalent, aus verschiedenen Perioden sind nicht äquivalent.}

Der erste Theil der Behauptung ist von vornherein klar, da reducirte Zahlen derselben Periode Schlusszahlen von einander sind; der zweite Theil ist durch den Satz \S.\,121, VI bewiesen, dass äquivalente Zahlen, bei hinlänglich weit fortgesetzter Ketten-
%CONT>>
